\documentclass[10pt,a4paper]{amsart}
\usepackage[margin=19mm]{geometry}
\usepackage{amsmath,amssymb,mathtools}
\usepackage[scr=rsfs,cal=euler]{mathalfa}
\usepackage{xfrac}
\usepackage[unicode,hidelinks,bookmarks=false]{hyperref}
\newcommand{\ie}{i.e.\ }
\newcommand{\cf}{cf.\ }
\newcommand{\resp}{respectively\ }
\newcommand{\Subsection}{\subsection}
\providecommand{\nobreakdash}{\nobreak\hbox{-}}
\setlength{\emergencystretch}{2em}
\multlinegap0pt
\usepackage{bm}
\usepackage{bbm}
\usepackage{dsfont}
\usepackage{xspace}
\usepackage{cancel}
\usepackage{mathtools}
\usepackage{amsthm}
\makeatletter
\@ifundefined{theorem}{\newtheorem{theorem}{Theorem}}{}
\makeatother
\newcommand{\N}{\mathbb{N}}
\newcommand{\Z}{\mathbb{Z}}
\newcommand{\abs}[1]{\left| #1 \right|}
\newcommand{\eps}{\varepsilon}
\newcommand{\eqNote}[1]{\stackrel{\text{#1}}{=}}
\newcommand{\leqNote}[1]{\stackrel{\text{#1}}{\leq}}
\newcommand{\norm}[1]{\left\| #1 \right\|}
\title[Dynamical properties of weighted shifts on sequence spaces]{Dynamical properties of weighted shifts on sequence spaces}
\author{Michal Hevessy, Tom\'a\v{s} Raunig}
\date{Source: arXiv:2509.11852v2 (CC BY 4.0)}
\begin{document}
\maketitle

\noindent\emph{Source and licence.} Dynamical properties of weighted shifts on sequence spaces. Version arXiv:2509.11852v2 (CC BY 4.0), distributed under the Creative Commons Attribution 4.0 International licence, as recorded in the arXiv licence metadata for this exact version and shown on the abstract page https://arxiv.org/abs/2509.11852v2. Full rights notice: licenses/arxiv-2509.11852-CC-BY-4.0.txt.\par\smallskip

\noindent\textit{Editorial scope.} Counted unit: the Theorem whose source label is \texttt{thm:PSPc0} in the article (source0.tex lines 745--756, source label \texttt{thm:PSPc0}), delivered together with the article's own complete original proof of it (source0.tex lines 757--802, one \texttt{proof} environment of 46 lines). No mathematics is written, repaired or re-derived: statement and proof are the article's own text line for line. Required context retained uncounted: none. Every definition, display and label of those lines is retained unchanged. Omitted from this package and not redistributed: 1--744, 803--813. Nothing inside a retained excerpt is omitted or altered: every retained body is delivered exactly as the article's lines are written, so no omission receipt of any kind accompanies this package. Citations: the retained excerpts carry no citation command at all, so no bibliography entry of the article is transcribed and no reference is redistributed. Reference adaptation: none was needed and none was made; every reference inside the delivered unit resolves inside it, so no unresolved reference or citation is produced. Printed slips kept literal: no printed word, symbol, hypothesis or number of the counted statement or of its proof was altered. Layout and class: the article's own class is \texttt{amsart} with packages including inputenc, babel, a4wide, amsmath, amssymb, amsthm, xcolor, csquotes, enumitem, hyperref, cleveref; the wrapper is the lane's pinned \texttt{amsart} editorial wrapper at 10pt on A4, which restates the theorem-like environments the retained text needs and adds the article's own macro definitions that the retained text needs (\texttt{\textbackslash N}, \texttt{\textbackslash Z}, \texttt{\textbackslash abs}, \texttt{\textbackslash eps}, \texttt{\textbackslash eqNote}, \texttt{\textbackslash leqNote}, \texttt{\textbackslash norm}), with extra wrapper packages (bm, bbm, dsfont, xspace, cancel, mathtools). The delivered numbering is the unit's own. Assets: no retained span contains a figure environment, a graphics-inclusion command, a PDF-page inclusion, a table or a TikZ picture; no image file is copied into this package, the package declares no asset dependency and no image file is redistributed with it. Licence: version arXiv:2509.11852v2 of this article is distributed under the Creative Commons Attribution 4.0 International licence, as recorded in the arXiv licence metadata for this exact version and shown on the abstract page https://arxiv.org/abs/2509.11852v2. A full rights notice is delivered at licenses/arxiv-2509.11852-CC-BY-4.0.txt. Characters: every retained excerpt is pure ASCII, so the delivered engine, XeLaTeX with its default Latin Modern text font, reports no missing character.
	\begin{theorem} \label{thm:PSPc0}
		Let $B_w$ be a~bilateral weighted backward shift on \(X = c_0(\Z)\) given by a~bounded sequence of weights \(w = (w_n)_{n \in \Z}\) with \(\inf_{n \in \Z}\abs{w_n} > 0\). If $B_w$ has nontrivial periodic points, then the periodic shadowing property is equivalent to the shadowing property. If not, the periodic shadowing property is equivalent to the condition:
		For every $\eps > 0$ there is $\delta > 0$ such that for all $k \leq l \leq m$ one of the following holds
		\begin{enumerate}
			\item[(a)] \begin{equation*}
				\abs{\frac{\varepsilon}{w(l+1)\dots w(m)}- \delta \sum_{i=0}^{m-(l+1)}\frac{1}{w(m) \dots w(m-i)}} \geq \delta
			\end{equation*}
			\item[(b)] \begin{equation*}
				\abs{w(k) \dots w(l)\varepsilon -\delta\sum_{i=0}^{k-(l-1)}w(k) \dots w(k-i)} \geq \delta.
			\end{equation*}
		\end{enumerate}
	\end{theorem}

	\begin{proof}
		We will show the reverse implication, proceeding by contraposition. Assume \(B_w\) does not have the PSP. Thus, there is \(\varepsilon > 0\) such that for every \(\eta > 0\) there is a \(\eta\)-pseudotrajectory for \(B_w\) starting at \(0\) such that it is not contained in \(B(0,\varepsilon)\).
		
		Pick \(\delta > 0\). Let \((z_i)_{i=0}^{n+2}\) be a periodic \(\delta\)-pseudotrajectory starting at \(0\) that is not contained in \(B(0,\varepsilon)\). For the ease of notation, define \(x_i = z_{i+1}\) for \(i \in \{0, \dots,n\}\). Then we have that \(\norm{x_0} < \delta \), \(\norm{B_w(x_{n})} < \delta \) and \((x_i)_{i=0}^{n}\) is a \(\delta\)-pseudotrajectory also not contained in \(B(0, \varepsilon)\). 
		
		Let \(r \in \N\) be such that \(\norm{x_r} \geq \varepsilon\). We can find \(l \in \Z\) such that \(\abs{x_r(l)} \geq \varepsilon\). Let \(y_0 = \eps e_l \) and for \(i \in \N\) define
		\begin{equation} \label{eqn:renormedErrorDef}
			y_{i}= \begin{cases}
				\left(1-\cfrac{\delta}{\norm{y_{i-1}}}\right)B_w^{-1}(y_{i-1}), & y_{i-1} \neq 0,\\
				0, & y_{i-1} = 0.
			\end{cases}
		\end{equation}
		This implies that every \(y_i\) has at most one non-zero coordinate.
		
		\textbf{Claim:} There is \(p \leq r\) such that \(\norm{y_p} < \delta\).
		
		Suppose this is not the case. Then we have \(\abs{y_r(l+r)} = \norm{y_r} \geq \delta > \abs{x_{0}(l+r)}\). Note that we have \(\abs{y_0(l)} \leq \abs{x_{r}(l)}\) by the definition of \(y_0\). Let \(u < r\) be maximal such that
		\begin{equation}\label{eqn:maximalu}
			\abs{y_u(l+u)} \leq \abs{x_{r-u}(l+u)}  
		\end{equation}
		(mind you, we have just shown that this holds for $u = 0$, but not $u=r$).
		
		By the assumption, we have that \(\delta \leq \abs{y_u(l+u)}\). This and the definition of \(y_i\) implies that
		\begin{equation} \label{eqn:stepIsDelta}
			\abs{B_w(y_{u+1})(l+u)}+\delta = \abs{y_{u}(l+u)}.    
		\end{equation}
		Since \((x_i)_{i=0}^{n}\) is a \(\delta\)-pseudotrajectory we have that
		\begin{equation} \label{eqn:stepIsLessThanDelta}
			\abs{(B_w(x_{r-u-1})-x_{r-u})(l+u)} < \delta.
		\end{equation}
		All together we get
		\begin{equation*}
			\abs{B_w(y_{u+1})(l+u)}
			\eqNote{\eqref{eqn:stepIsDelta}} \abs{y_u(l+u)}-\delta
			\leqNote{\eqref{eqn:maximalu}} \abs{x_{r-u}(l+u)}-\delta
			\leqNote{\eqref{eqn:stepIsLessThanDelta}} \abs{B_w(x_{r-u-1})(l+u)}.
		\end{equation*}
		But since \(B^{-1}_w\) preserves the pointwise order we get that \(\abs{y_{u+1}(l+u+1)} \leq \abs{x_{r-u-1}(l+u+1)}\), which contradicts the maximality of \(u\). Thus the claim is proved.
		
		
		Let \(p\) be such that \(\norm{y_p} < \delta\). Using simple induction and~\eqref{eqn:renormedErrorDef}, one can show that
		\begin{align*}
			\delta > \norm{y_p} = \abs{y_p(l+p)} = \abs{\frac{\varepsilon}{w(l+1)\dots w(l+p)}- \delta \sum_{i=0}^{1-p}\frac{1}{w(l+p) \dots w(l+p+j)}}. 
		\end{align*}
		By taking \(l+p = m\) we get the first part of the claim. The second part can be shown similarly by starting with \(z_0 = B_w(y_0)\) and doing similar construction using \(B_w\) instead of \(B^{-1}_w\).
	\end{proof}

\end{document}
